\section{Discussion}
\label{sec:discussion}
\subsection{What limits the magnification of a lensed GW image?}
Our results separate five sources of magnification error. (i) \emph{Lens-model freedom at fixed image configuration.} This dominates for realistic lenses: from SIS to SIE, SIE+shear and EPL+shear the geometry-only error grows from 3 to 7, 12 and 20--21\% for ET. The radial density slope is the most important unobserved parameter, as expected from the slope--magnification degeneracy familiar from time-delay cosmography. (ii) \emph{Proximity to a caustic.} In our populations the errors grow roughly in proportion to $\mu$, so the most strongly magnified---and, for current detectors, most likely to be detected---images are the hardest to correct. (iii) \emph{The GW apparent distance.} Under the catalog-calibrated error model, the inclination--distance degeneracy limits each image's apparent distance to $\sigma_{\ln D}\approx0.22$, and this dominates the error of the corrected luminosity distance. Joint PE of both images with higher harmonics gives apparent-distance errors smaller by a factor 2.3 for ET and 1.5 for O5 (Sec.~\ref{sec:jointpe}), so this conclusion depends on the GW analysis and should not be read as a fundamental floor. (iv) \emph{A wrong lens prior.} An estimator trained on an over-simplified lens population produces biased estimates with underestimated uncertainties. This is the most dangerous failure mode, because a narrow interval by itself does not reveal an inadequate mass model; residual and model-comparison diagnostics are needed to detect it. (v) \emph{An implicit distance scale.} An estimator that receives the delay or the SNRs together with redshifts gains accuracy by learning the simulated cosmology, and its output then shifts with the assumed $H_0$ (Table~\ref{tab:response}).

These findings suggest a practical strategy for the first lensed BBH with an identified host. Image geometry and the GW relative magnification should be combined with a lens prior that is at least as flexible as an EPL plus external shear with an empirically calibrated slope distribution. Intervals should be calibrated on that flexible prior, because calibration groups defined by observables do not reproduce per-family calibration (Sec.~\ref{sec:calibration}). Inputs that carry a distance scale should be used only when that scale is part of the analysis rather than an unexamined assumption. Any additional handle on the radial slope---stellar kinematics with a dynamical model, extended-arc modeling of the host, or the GW time delays themselves combined with an independent time-delay distance---should be added, because it targets the dominant error term.

\subsection{Mass-sheet degeneracy and cosmology dependence}
Every result above is conditional on the simulated lens populations, including their distribution of radial profiles, and on the absence of an external mass sheet. For estimators that use only scale-free inputs, a sheet with $\lambda\neq1$ multiplies the estimate relative to the truth by exactly $\lambda^2$, which we verified numerically (Table~\ref{tab:response}). For a typical external convergence $|\kappa_{\rm ext}|\lesssim0.05$, this bias of up to 10\% in $\mu$ (5\% in $d_L$) is comparable to our SIE errors and below those of the EPL family. It would, however, be coherent across a population of lensed events if environments were systematically over- or under-dense. Estimators that use the delay respond to $\lambda$, but only partially (Table~\ref{tab:response}), so they do not remove the bias either. Line-of-sight convergence estimates from galaxy counts or ray tracing, as in time-delay cosmography, can be folded into the prior. Any constraint that uses an assumed cosmology, however, makes the corrected distance circular for cosmography, so the scale-free results are the relevant benchmark for that application.

\subsection{Limitations}
We assumed that the GW images have been associated with each other and with the EM images of an identified host, and that the image positions of the BBH within the host are reconstructed to the stated astrometric precision. The BBH itself is not observed in the EM data, so these positions are the result of a preceding reconstruction of the extended host and the GW information. Its uncertainty may be larger, correlated between images and multimodal, and it must not reuse the GW measurements that enter the magnification inference. Repeating the Bayesian reconstruction of Sec.~\ref{sec:gwem} with positional errors of 0.1\,arcsec raises the errors from 7.9 to 12.0\% (ET) and from 13.3 to 14.4\% (O5), with coverage 0.92 and 0.86, provided the exact position likelihood of Appendix~\ref{app:gwemlike} is used; the linearized one gives coverage of only 0.70 and 0.64. A pixel-level reconstruction of lensed hosts with off-center BBHs, propagated into the GW-source image positions, is needed to replace this assumption; we leave it to future work. The association steps are probabilistic, and the magnification posterior must marginalize over candidate associations and lens reconstructions \citep{hannuksela2020localizing,wempe2022lensing,uronen2026galacticcentres}. The benchmark itself uses the $(2,2)$ mode of an aligned-spin waveform, an SNR-based amplitude model and analytic scalings for the sky and orientation errors. The O4a injections, which recover each image with the exact template, test the relative-amplitude noise prescription in the LIGO detectors only; the joint PE of Sec.~\ref{sec:jointpe} tests the SNR-based measurement model with higher harmonics and free source parameters, but in Gaussian noise and with aligned spins. The conventional Bayesian reconstructions assume the simulated lens family unless stated otherwise, and their nominal 90\% intervals contain the truth for only about 80\% of the ET systems (Sec.~\ref{sec:gwem}). Lens populations do not include satellites, line-of-sight structures or redshift evolution of lens properties, and wave-optics effects are negligible for the galaxy-scale lenses considered here but would matter for lower-mass lenses. Each of these simplifications should be relaxed before the method is applied to data.

\section{Conclusions}
\label{sec:conclusions}
We have benchmarked the accuracy and reliability of absolute-magnification inference for strongly lensed GW images from GW and EM observables. The benchmark uses GWTC-4.0 source populations, four galaxy-lens families, O5-like and ET networks, GW distance errors calibrated on 167 GWTC events, and tests of the GW measurement model in public O4a noise and by joint parameter estimation. Our main conclusions are as follows.
\begin{enumerate}
\item Image geometry measured relative to the lens-light center is the most informative observable. Combined with the GW magnification ratio and arrival order---observables that are invariant under an external mass sheet and contain no distance scale---it determines the magnification of an ET image to a median 3.0, 6.9, 10.2 and 18.8\% for SIS, SIE, SIE+shear and EPL+shear lenses (HistGB trained on the correct family; NPE trained on all families gives 3.3--19.6\%). For the more magnification-biased O5 sample the errors are 10.8--25.5\%. GW data alone give 21--26\% for realistic lenses.
\item The time delay and the SNRs add little on their own. With redshifts they reduce the errors further, to 1.7--13.6\% for NPE, but only by learning the distance scale of the training cosmology: the inferred magnification then shifts by 0.7\% per percent change in $H_0$. Such estimates are not appropriate for cosmography.
\item Within a lens family, the calibrated (conformal) 90\% intervals contain the truth 90\% of the time. An estimator trained on all families needs the family label to be calibrated per family; groups of events defined by observables do not substitute for it. Near caustics, the intervals contain the truth too rarely for the truly highly magnified images (44--57\% for $\mu_A>20$).
\item The lens prior is the dominant systematic. SIS-trained models give biased estimates with underestimated uncertainties for realistic lenses (coverage 20--27\%), and omitting radial-slope freedom drops the coverage to 55\%. Calibrating on the flexible EPL+shear population gives coverage of 0.89--0.99 in every family at 40\% wider intervals.
\item The SNR-based GW measurement model is conservative. It holds in real O4a noise for exact-template recovery (standardized scatter 1.02), and joint parameter estimation of 120 image pairs from simulated detector data, with higher harmonics and free source parameters, gives smaller uncertainties, by a factor 2.3 for the apparent distance of ET images (1.5 for O5). The apparent distance therefore dominates the error of the corrected luminosity distance only under the catalog-calibrated error model (22--24\% total for ET, of which 1.5--9.5\% comes from the magnification). For ET the host sky position is needed: with a free sky position the posterior favors a wrong sky position in a quarter of the systems.
\item A conventional Bayesian GW+EM reconstruction with the same observables reaches 7.9 (ET) and 13.7\% (O5) when it knows the lens family, comparable to the NPE (10.5 and 13.8\%), and 12.4 and 18.6\% with an EPL+shear model for every system, provided the prior of the source position is normalized for each lens as in the simulated population. Its intervals are 40--45\% narrower, but for ET they contain the truth for only about 80\% of the systems.
\end{enumerate}
The code, configuration and real-data calibration are public. The results are conditional on reconstructed GW-source image positions and on the absence of an external mass sheet; pixel-level reconstruction of lensed hosts with off-center BBHs and line-of-sight convergence are the most important next steps.
